 Capítulo  7
 Agendamento
 Qualquer  sistema  operacional  provavelmente  executará  com  mais  processos  do  que  a  capacidade  de  CPU  do  
computador,  portanto,  é  necessário  um  plano  para  compartilhar  o  tempo  das  CPUs  entre  os  processos.  Idealmente,  o  
compartilhamento  seria  transparente  para  os  processos  do  usuário.  Uma  abordagem  comum  é  dar  a  cada  processo  a  
ilusão  de  que  possui  sua  própria  CPU  virtual,  multiplexando  os  processos  nas  CPUs  de  hardware.  Este  capítulo  explica  
como  o  xv6  realiza  essa  multiplexação.
 7.1  Multiplexação
 O  Xv6  multiplexa  alternando  cada  CPU  de  um  processo  para  outro  em  duas  situações.  Primeiro,  o  mecanismo  de  
suspensão  e  ativação  do  xv6  alterna  quando  um  processo  aguarda  a  conclusão  da  E/S  do  dispositivo  ou  do  pipe,  ou  
aguarda  a  saída  de  um  processo  filho,  ou  aguarda  na  chamada  de  sistema  de  suspensão .  Segundo,  o  xv6  força  
periodicamente  uma  troca  para  lidar  com  processos  que  computam  por  longos  períodos  sem  suspensão.  Essa  
multiplexação  cria  a  ilusão  de  que  cada  processo  possui  sua  própria  CPU,  assim  como  o  xv6  utiliza  o  alocador  de  
memória  e  as  tabelas  de  páginas  de  hardware  para  criar  a  ilusão  de  que  cada  processo  possui  sua  própria  memória.
 A  implementação  da  multiplexação  apresenta  alguns  desafios.  Primeiro,  como  alternar  de  um  processo  para  
outro?  Embora  a  ideia  de  troca  de  contexto  seja  simples,  a  implementação  é  um  dos  códigos  mais  opacos  do  Xv6.  
Segundo,  como  forçar  trocas  de  forma  transparente  para  os  processos  do  usuário?  O  Xv6  usa  a  técnica  padrão  na  
qual  as  interrupções  de  um  temporizador  de  hardware  controlam  as  trocas  de  contexto.  Terceiro,  todas  as  CPUs  
alternam  entre  o  mesmo  conjunto  compartilhado  de  processos,  e  um  plano  de  bloqueio  é  necessário  para  evitar  
corridas.  Quarto,  a  memória  e  outros  recursos  de  um  processo  devem  ser  liberados  quando  o  processo  termina,  mas  
ele  não  pode  fazer  tudo  isso  sozinho  porque  (por  exemplo)  não  pode  liberar  sua  própria  pilha  do  kernel  enquanto  ainda  
a  utiliza.  Quinto,  cada  núcleo  de  uma  máquina  multinúcleo  deve  se  lembrar  de  qual  processo  está  executando  para  
que  as  chamadas  de  sistema  afetem  o  estado  correto  do  kernel  do  processo.  Finalmente,  a  suspensão  e  a  ativação  
permitem  que  um  processo  desista  da  CPU  e  espere  ser  ativado  por  outro  processo  ou  interrupção.  É  necessário  
cuidado  para  evitar  corridas  que  resultem  na  perda  de  notificações  de  ativação.  O  Xv6  tenta  resolver  esses  problemas  
da  forma  mais  simples  possível,  mas  mesmo  assim  o  código  resultante  é  complicado.
 71
Machine Translated by Google
 usuário
 espaço
 núcleo
 espaço
 concha
 salvar
 kstack
 concha
 gato
 interruptor
 interruptor
 agendador  
kstack
 Núcleo
 kstack
 gato
 restaurar
 Figura  7.1:  Troca  de  um  processo  de  usuário  para  outro.  Neste  exemplo,  o  xv6  é  executado  com  uma  CPU  (e,  portanto,  
uma  thread  do  escalonador).
 7.2  Código:  Troca  de  contexto
 A  Figura  7.1  descreve  as  etapas  envolvidas  na  alternância  de  um  processo  de  usuário  para  outro:  uma  transição  usuário
kernel  (chamada  de  sistema  ou  interrupção)  para  a  thread  do  kernel  do  processo  antigo,  uma  troca  de  contexto  para  a  
thread  do  escalonador  da  CPU  atual,  uma  troca  de  contexto  para  a  thread  do  kernel  de  um  novo  processo  e  um  retorno  
de  trap  para  o  processo  em  nível  de  usuário.  O  escalonador  xv6  possui  uma  thread  dedicada  (registradores  e  pilha  salvos)  
por  CPU,  pois  não  é  seguro  para  o  escalonador  executar  na  pilha  do  kernel  do  processo  antigo:  algum  outro  núcleo  pode  
despertar  o  processo  e  executá-lo,  e  seria  um  desastre  usar  a  mesma  pilha  em  dois  núcleos  diferentes.  Nesta  seção,  
examinaremos  a  mecânica  da  alternância  entre  uma  thread  do  kernel  e  uma  thread  do  escalonador.
 Alternar  de  um  thread  para  outro  envolve  salvar  os  registradores  da  CPU  do  thread  antigo  e  restaurar  os  registradores  
salvos  anteriormente  do  novo  thread;  o  fato  de  o  ponteiro  de  pilha  e  o  contador  de  programa  serem  salvos  e  restaurados  
significa  que  a  CPU  alternará  as  pilhas  e  o  código  que  está  executando.
 A  função  swtch  executa  salvamentos  e  restaurações  para  uma  troca  de  thread  do  kernel.  swtch  não  sabe  diretamente  
sobre  threads;  ele  apenas  salva  e  restaura  conjuntos  de  32  registradores  RISC-V,  chamados  contextos.
 Quando  chega  a  hora  de  um  processo  liberar  a  CPU,  a  thread  do  kernel  do  processo  chama  swtch  para  salvar  seu  próprio  
contexto  e  retornar  ao  contexto  do  escalonador.  Cada  contexto  está  contido  em  uma  estrutura  context  (kernel/proc.h:2).  
contido  na  struct  proc  de  um  processo  ou  na  struct  cpu  de  uma  CPU .  swtch  recebe  dois  argumentos:  struct  context  *old  e  
struct  context  *new.  Ele  salva  os  registradores  atuais  em  old,  carrega  os  registradores  de  new  e  retorna.
 Vamos  acompanhar  um  processo  através  de  swtch  até  o  escalonador.  Vimos  no  Capítulo  4  que  uma  
possibilidade  ao  final  de  uma  interrupção  é  que  a  usertrap  chame  yield.  O  yield,  por  sua  vez,  chama  sched,  
que  chama  swtch  para  salvar  o  contexto  atual  em  p->context  e  alternar  para  o  contexto  do  escalonador  salvo  
anteriormente  em  cpu->context  (kernel/proc.c:497).
 interruptor  (kernel/swtch.S:3)  Salva  apenas  os  registradores  salvos  pelo  chamador;  o  compilador  C  gera  código  no  
chamador  para  salvar  os  registradores  salvos  pelo  chamador  na  pilha.  O  swtch  sabe  o  deslocamento  do  campo  de  cada  
registrador  no  contexto  da  estrutura.  Ele  não  salva  o  contador  do  programa.  Em  vez  disso,  o  swtch  salva  o  registrador  ra ,  
que  contém  o  endereço  de  retorno  de  onde  o  swtch  foi  chamado.  Agora,  o  swtch  restaura  os  registradores  de
 72
Machine Translated by Google
 O  novo  contexto,  que  contém  valores  de  registradores  salvos  por  um  swtch  anterior.  Quando  swtch  retorna,  ele  retorna  
às  instruções  apontadas  pelo  registrador  ra  restaurado ,  ou  seja,  a  instrução  da  qual  a  nova  thread  chamou  swtch  
anteriormente.  Além  disso,  ele  retorna  para  a  pilha  da  nova  thread,  pois  é  para  lá  que  o  sp  restaurado  aponta.
 Em  nosso  exemplo,  o  sched  chamou  swtch  para  alternar  para  cpu->context,  o  contexto  do  escalonador  por  CPU.  
Esse  contexto  foi  salvo  no  momento  em  que  o  escalonador  chamou  swtch  (kernel /proc.c:463).  para  alternar  para  o  
processo  que  está  desistindo  da  CPU.  Quando  o  switch  que  estávamos  rastreando  retorna,  ele  não  retorna  para  sched ,  
mas  para  scheduler,  com  o  ponteiro  de  pilha  na  pilha  do  scheduler  da  CPU  atual.
 7.3  Código:  Agendamento
 A  última  seção  analisou  os  detalhes  de  baixo  nível  do  swtch;  agora,  vamos  considerar  o  swtch  como  um  dado  e  examinar  
a  comutação  de  uma  thread  do  kernel  de  um  processo  para  outro,  através  do  escalonador.  O  escalonador  existe  na  
forma  de  uma  thread  especial  por  CPU,  cada  uma  executando  a  função  de  escalonador .  Esta  função  é  responsável  
por  escolher  qual  processo  executar  em  seguida.  Um  processo  que  deseja  abrir  mão  da  CPU  deve  adquirir  seu  próprio  
bloqueio  de  processo  p->lock,  liberar  quaisquer  outros  bloqueios  que  esteja  mantendo,  atualizar  seu  próprio  estado  (p
>state)  e,  em  seguida,  chamar  sched.  Você  pode  ver  esta  sequência  em  yield  (kernel/proc.c:503).  dormir  e  sair.  sched  
verifica  novamente  alguns  desses  requisitos  (kernel/proc.c:487-492)  e  então  verifica  uma  implicação:  como  um  bloqueio  
é  mantido,  as  interrupções  devem  ser  desabilitadas.  Por  fim,  sched  chama  swtch  para  salvar  o  contexto  atual  em  p
>context  e  alternar  para  o  contexto  do  escalonador  em  cpu->context.  swtch  retorna  na  pilha  do  escalonador  como  se  o  
swtch  do  escalonador  tivesse  retornado  (kernel/proc.c:463).  O  agendador  continua  seu  loop  for ,  encontra  um  processo  
para  executar,  alterna  para  ele  e  o  ciclo  se  repete.
 Acabamos  de  ver  que  o  xv6  mantém  p->lock  em  chamadas  para  swtch:  o  chamador  de  swtch  já  deve  manter  o  
bloqueio,  e  o  controle  do  bloqueio  passa  para  o  código  comutado.  Essa  convenção  é  incomum  com  bloqueios;  
geralmente,  a  thread  que  adquire  um  bloqueio  também  é  responsável  por  liberá-lo,  o  que  facilita  o  raciocínio  sobre  a  
correção.  Para  troca  de  contexto,  é  necessário  quebrar  essa  convenção  porque  p->lock  protege  invariantes  nos  campos  
de  estado  e  contexto  do  processo  que  não  são  verdadeiros  durante  a  execução  em  swtch.  Um  exemplo  de  um  problema  
que  poderia  surgir  se  p->lock  não  fosse  mantido  durante  swtch:  uma  CPU  diferente  poderia  decidir  executar  o  processo  
após  yield  ter  definido  seu  estado  como  RUNNABLE,  mas  antes  de  swtch  fazer  com  que  ele  parasse  de  usar  sua  
própria  pilha  do  kernel.  O  resultado  seriam  duas  CPUs  rodando  na  mesma  pilha,  o  que  causaria  o  caos.
 O  único  lugar  onde  uma  thread  do  kernel  cede  sua  CPU  é  no  sched,  e  ela  sempre  alterna  para  o  mesmo  local  no  
escalonador,  que  (quase)  sempre  alterna  para  alguma  thread  do  kernel  que  anteriormente  chamava  o  sched.  Portanto,  
se  alguém  imprimisse  os  números  das  linhas  onde  o  xv6  alterna  as  threads,  observaria  o  seguinte  padrão  simples:  
(kernel/proc.c:463),  (kernel/proc.c:497),  (kernel/proc.c:463),  (kernel/proc.c:497),  e  assim  por  diante.  Procedimentos  que  
intencionalmente  transferem  o  controle  entre  si  por  meio  de  troca  de  threads  são  às  vezes  chamados  de  corrotinas;  
neste  exemplo,  sched  e  scheduler  são  corrotinas  uma  da  outra.
 Há  um  caso  em  que  a  chamada  do  planejador  para  swtch  não  termina  em  sched.  allocproc  define  o  registro  ra  de  
contexto  de  um  novo  processo  para  forkret  (kernel/proc.c:515),  para  que  sua  primeira  troca
 73
Machine Translated by Google
 “retorna”  ao  início  dessa  função.  forkret  existe  para  liberar  o  p->lock;  caso  contrário,  como  o  novo  processo  precisa  retornar  
ao  espaço  do  usuário  como  se  estivesse  retornando  de  fork,  ele  poderia  começar  em  usertrapret.  scheduler  (kernel/
 proc.c:445)  executa  
um  loop:  encontre  um  processo  para  executar,  execute-o  até  que  ele  ceda  e  repita.
 O  escalonador  percorre  a  tabela  de  processos  em  busca  de  um  processo  executável,  que  tenha  p->state  ==  
RUNNABLE.  Ao  encontrar  um  processo,  ele  define  a  variável  de  processo  atual  por  CPU  c->proc,  marca  o  
processo  como  EM  EXECUÇÃO  e,  em  seguida,  chama  swtch  para  iniciar  sua  execução  (kernel/proc.c:458-463).
 Uma  maneira  de  pensar  sobre  a  estrutura  do  código  de  escalonamento  é  que  ele  impõe  um  conjunto  de  invariantes  
sobre  cada  processo  e  mantém  p->lock  sempre  que  essas  invariantes  não  forem  verdadeiras.  Uma  invariante  é  que,  se  um  
processo  estiver  EM  EXECUÇÃO,  o  yield  de  uma  interrupção  de  temporizador  deve  ser  capaz  de  alternar  com  segurança  
para  longe  do  processo;  isso  significa  que  os  registradores  da  CPU  devem  conter  os  valores  de  registrador  do  processo  (ou  
seja,  swtch  não  os  moveu  para  um  contexto)  e  c->proc  deve  se  referir  ao  processo.  Outra  invariante  é  que,  se  um  processo  
for  EXECUTÁVEL,  deve  ser  seguro  para  o  escalonador  de  uma  CPU  ociosa  executá-lo;  isso  significa  que  p->context  deve  
conter  os  registradores  do  processo  (ou  seja,  eles  não  estão  realmente  nos  registradores  reais),  que  nenhuma  CPU  está  
executando  na  pilha  do  kernel  do  processo  e  que  nenhum  c->proc  da  CPU  se  refere  ao  processo.  Observe  que  essas  
propriedades  geralmente  não  são  verdadeiras  enquanto  p->lock  é  mantido.
 Manter  as  invariantes  acima  é  o  motivo  pelo  qual  o  xv6  frequentemente  adquire  p->lock  em  uma  thread  e  o  libera  em  
outra,  por  exemplo,  adquirindo  em  yield  e  liberando  em  scheduler.  Uma  vez  que  yield  tenha  começado  a  modificar  o  estado  
de  um  processo  em  execução  para  torná-lo  RUNNABLE,  o  bloqueio  deve  permanecer  mantido  até  que  as  invariantes  sejam  
restauradas:  o  primeiro  ponto  de  liberação  correto  é  após  o  scheduler  (executando  em  sua  própria  pilha)  limpar  c->proc.  Da  
mesma  forma,  uma  vez  que  o  scheduler  começa  a  converter  um  processo  RUNNABLE  para  RUNNING,  o  bloqueio  não  
pode  ser  liberado  até  que  a  thread  do  kernel  esteja  completamente  em  execução  (após  a  troca,  por  exemplo  em  yield).
 7.4  Código:  mycpu  e  myproc
 O  Xv6  frequentemente  precisa  de  um  ponteiro  para  a  estrutura  proc  do  processo  atual .  Em  um  processador  único,  pode-se  
ter  uma  variável  global  apontando  para  o  proc  atual.  Isso  não  funciona  em  uma  máquina  com  vários  núcleos,  já  que  cada  
núcleo  executa  um  processo  diferente.  A  maneira  de  resolver  esse  problema  é  explorar  o  fato  de  que  cada  núcleo  tem  seu  
próprio  conjunto  de  registradores;  podemos  usar  um  desses  registradores  para  ajudar  a  encontrar  informações  por  núcleo.
 O  Xv6  mantém  uma  estrutura  cpu  para  cada  CPU  (kernel/proc.h:22),  que  registra  o  processo  em  execução  naquela  
CPU  (se  houver),  os  registradores  salvos  para  a  thread  do  escalonador  da  CPU  e  a  contagem  de  spinlocks  aninhados  
necessários  para  gerenciar  a  desabilitação  de  interrupções.  A  função  mycpu  (kernel/proc.c:74)  O  Xv6  retorna  um  ponteiro  
para  a  estrutura  cpu  da  CPU  atual.  O  RISC-V  numera  suas  CPUs,  atribuindo  a  cada  uma  um  hartid.  O  Xv6  garante  que  o  
hartid  de  cada  CPU  seja  armazenado  no  registrador  tp  dessa  CPU  enquanto  estiver  no  kernel.
 Isso  permite  que  o  mycpu  use  o  tp  para  indexar  uma  matriz  de  estruturas  de  CPU  para  encontrar  a  correta.
 Garantir  que  o  tp  de  uma  CPU  sempre  mantenha  o  hartid  da  CPU  é  um  pouco  complicado.  start  define  o  
registrador  tp  no  início  da  sequência  de  inicialização  da  CPU,  enquanto  ainda  está  no  modo  de  máquina  (kernel/
 start.c:51).  usertrapret  salva  tp  na  página  do  trampolim,  pois  o  processo  do  usuário  pode  modificá  -lo.  Por  fim,  
uservec  restaura  esse  tp  salvo  ao  entrar  no  kernel  a  partir  do  espaço  do  usuário  (kernel/trampoline.S:77).
 O  compilador  garante  nunca  usar  o  registrador  tp .  Seria  mais  conveniente  se  o  xv6  pudesse
 74
Machine Translated by Google
 pergunte  ao  hardware  RISC-V  sobre  o  hartid  atual  sempre  que  necessário,  mas  o  RISC-V  permite  isso  apenas  em
 modo  máquina,  não  no  modo  supervisor.
 Os  valores  de  retorno  de  cpuid  e  mycpu  são  frágeis:  se  o  temporizador  interrompesse  e  causasse  o
 thread  para  render  e  então  mover  para  uma  CPU  diferente,  um  valor  retornado  anteriormente  não  seria  mais
 correto.  Para  evitar  esse  problema,  o  xv6  requer  que  os  chamadores  desabilitem  as  interrupções  e  as  habilitem  apenas
 depois  que  eles  terminam  de  usar  a  estrutura  retornada  cpu.
 A  função  myproc  (kernel/proc.c:83)  retorna  o  ponteiro  struct  proc  para  o  processo  que
 está  sendo  executado  na  CPU  atual.  myproc  desabilita  interrupções,  invoca  mycpu,  busca  a  CPU  atual
 ponteiro  de  processo  (c->proc)  para  fora  da  estrutura  da  CPU  e,  em  seguida,  habilita  interrupções.  O  valor  de  retorno
 de  myproc  é  seguro  de  usar  mesmo  se  as  interrupções  estiverem  habilitadas:  se  uma  interrupção  do  temporizador  move  o  processo  de  chamada
 para  uma  CPU  diferente,  seu  ponteiro  struct  proc  permanecerá  o  mesmo.
 7.5  Sono  e  despertar
 Agendamentos  e  bloqueios  ajudam  a  ocultar  as  ações  de  uma  thread  de  outra,  mas  também  precisamos  de  abstrações  que  
ajudem  as  threads  a  interagir  intencionalmente.  Por  exemplo,  o  leitor  de  um  pipe  em  xv6  pode
 precisa  esperar  que  um  processo  de  escrita  produza  dados;  a  chamada  de  um  pai  para  esperar  pode  precisar  esperar  por  um
 criança  para  sair;  e  um  processo  que  lê  o  disco  precisa  esperar  que  o  hardware  do  disco  conclua  a  leitura.
 O  kernel  xv6  usa  um  mecanismo  chamado  suspensão  e  ativação  nessas  situações  (e  em  muitas  outras).
 O  modo  de  suspensão  permite  que  uma  thread  do  kernel  aguarde  um  evento  específico;  outra  thread  pode  chamar  o  modo  de  ativação  para  indicar  
que  as  threads  que  aguardam  um  evento  devem  ser  retomadas.  O  modo  de  suspensão  e  a  ativação  são  frequentemente  chamados  de  sequências.
 mecanismos  de  coordenação  ou  sincronização  condicional.
 O  sono  e  a  vigília  fornecem  uma  interface  de  sincronização  de  nível  relativamente  baixo.  Para  motivar  o
 como  eles  funcionam  no  xv6,  nós  os  usaremos  para  construir  um  mecanismo  de  sincronização  de  nível  superior  chamado
 um  semáforo  [5]  que  coordena  produtores  e  consumidores  (xv6  não  usa  semáforos).  A
 O  semáforo  mantém  uma  contagem  e  fornece  duas  operações.  A  operação  “V”  (para  o  produtor)
 incrementa  a  contagem.  A  operação  “P”  (para  o  consumidor)  espera  até  que  a  contagem  seja  diferente  de  zero,
 e  então  decrementa  e  retorna.  Se  houvesse  apenas  um  thread  produtor  e  um  consumidor
 thread,  e  eles  foram  executados  em  CPUs  diferentes,  e  o  compilador  não  otimizou  de  forma  muito  agressiva,
 esta  implementação  estaria  correta:
 100  estrutura  semáforo  {
 101  estrutura  de  bloqueio  spinlock;
 102
 103
 104
 105  vazio
 contagem  de  int;
 };
 106  V(estrutura  semáforo  *s)
 107  {
 108
 109
 110
 111
 adquirir(&s->bloquear);
 s->contagem  +=  1;
 liberar(&s->bloquear);
 }
 75
Machine Translated by Google
 112
 113  vazio
 114  P(struct  semáforo  *s)
 115  {
 116
 117
 118  
119  
120
 121
 enquanto(s->contagem  ==  0)
 ;
 adquirir(&s->bloquear);
 s->contagem  -=  1;
 liberar(&s->bloquear);
 }
 A  implementação  acima  é  cara.  Se  o  produtor  agir  raramente,  o  consumidor  gastará
 a  maior  parte  do  tempo  girando  no  loop  while  esperando  por  uma  contagem  diferente  de  zero.  A  CPU  do  consumidor
 provavelmente  conseguiria  encontrar  trabalho  mais  produtivo  do  que  ficar  ocupado  esperando,  consultando  repetidamente  s->count.
 Evitar  a  espera  ocupada  requer  uma  maneira  para  o  consumidor  ceder  a  CPU  e  retomar  somente  após  V
 incrementa  a  contagem.
 Aqui  está  um  passo  nessa  direção,  embora,  como  veremos,  não  seja  suficiente.  Imaginemos  um  par
 de  chamadas,  sleep  e  wakeup,  que  funcionam  da  seguinte  maneira.  sleep(chan)  dorme  no  valor  arbitrário
 chan,  chamado  de  canal  de  espera.  sleep  coloca  o  processo  de  chamada  em  hibernação,  liberando  a  CPU  para  outros
 trabalho.  wakeup(chan)  desperta  todos  os  processos  que  estão  dormindo  no  chan  (se  houver),  fazendo  com  que  suas  chamadas  de  sono  sejam
 retornar.  Se  não  houver  processos  aguardando  o  canal,  o  wakeup  não  faz  nada.  Podemos  alterar  o  semáforo
 implementação  para  usar  sleep  e  wakeup  (mudanças  destacadas  em  amarelo):
 200  vazio
 201  V(struct  semáforo  *s)
 202  {
 203
 204
 205
 206  
207
 208
 209  nulo
 adquirir(&s->bloquear);
 s->contagem  +=  1;
 despertar(es);
 liberar(&s->bloquear);
 }
 210  P(estrutura  semáforo  *s)
 211  {
 212
 213
 214
 215  
216
 217
 enquanto(s->contagem  ==  0)
 sono(s);
 adquirir(&s->bloquear);
 s->contagem  -=  1;
 liberar(&s->bloquear);
 }
 P  agora  desiste  da  CPU  em  vez  de  girar,  o  que  é  bom.  No  entanto,  não  é
 fácil  de  projetar  o  sono  e  o  despertar  com  esta  interface  sem  sofrer  com  o  que  é
 conhecido  como  o  problema  do  despertar  perdido.  Suponha  que  P  encontre  s->count  ==  0  na  linha  212.  Enquanto
 P  está  entre  as  linhas  212  e  213,  V  roda  em  outra  CPU:  ele  muda  s->count  para  ser  diferente  de  zero  e
 76
Machine Translated by Google
 chama  wakeup,  que  não  encontra  nenhum  processo  dormindo  e,  portanto,  não  faz  nada.  Agora  P  continua  executando
 na  linha  213:  ele  chama  sleep  e  entra  em  modo  sleep.  Isso  causa  um  problema:  P  está  dormindo  esperando  uma  chamada  de  V
 isso  já  aconteceu.  A  menos  que  tenhamos  sorte  e  o  produtor  ligue  para  V  novamente,  o  consumidor  irá
 esperar  para  sempre,  mesmo  que  a  contagem  seja  diferente  de  zero.
 A  raiz  deste  problema  é  que  a  invariante  de  que  P  só  dorme  quando  s->count  ==  0  é  violada
 por  V  correndo  no  momento  errado.  Uma  maneira  incorreta  de  proteger  a  invariante  seria  mover
 a  aquisição  do  bloqueio  (destacada  em  amarelo  abaixo)  em  P  para  que  sua  verificação  da  contagem  e  sua  chamada  para
 o  sono  é  atômico:
 300  vazio
 301  V(struct  semáforo  *s)
 302  {
 303
 304
 305
 306
 307
 308
 309  nulo
 adquirir(&s->bloquear);
 s->contagem  +=  1;
 despertar(es);
 liberar(&s->bloquear);
 }
 310  P(estrutura  semáforo  *s)
 311  {
 312
 313  
314  
315  
316
 317
 adquirir(&s->bloquear);
 enquanto(s->contagem  ==  0)
 sono(s);
 s->contagem  -=  1;
 liberar(&s->bloquear);
 }
 Poderíamos  esperar  que  esta  versão  de  P  evitasse  o  despertar  perdido  porque  o  bloqueio  impede  V
 de  executar  entre  as  linhas  313  e  314.  Ele  faz  isso,  mas  também  gera  um  deadlock:  P  mantém  o  bloqueio
 enquanto  ele  dorme,  então  V  ficará  bloqueado  para  sempre  esperando  o  bloqueio.
 Vamos  corrigir  o  esquema  anterior  alterando  a  interface  do  sleep :  o  chamador  deve  passar  o  bloqueio  de  condição  para  sleep  para  
que  ele  possa  liberar  o  bloqueio  depois  que  o  processo  de  chamada  for  marcado  como  sleep  e
 aguardando  no  canal  de  suspensão.  O  bloqueio  forçará  um  V  concorrente  a  esperar  até  que  P  termine  de  colocar
 para  dormir,  para  que  o  despertador  encontre  o  consumidor  adormecido  e  o  acorde.  Assim  que  o  consumidor  acorda  novamente,  o  sono  
recupera  o  bloqueio  antes  de  retornar.  Nosso  novo  modo  correto  de  sono/despertar
 o  esquema  pode  ser  utilizado  da  seguinte  forma  (alteração  destacada  em  amarelo):
 400  vazio
 401  V(struct  semáforo  *s)
 402  {
 403
 404
 405
 406
 407
 adquirir(&s->bloquear);
 s->contagem  +=  1;
 despertar(es);
 liberar(&s->bloquear);
 }
 77
Machine Translated by Google
 408
 409  nulo
 410  P(struct  semáforo  *s)  411  {
 412
 413
 414
 415
 416
 417
 adquirir(&s->bloquear);  
enquanto(s->contagem  ==  0)
 dormir(s,  &s->bloqueio);  s
>contagem  -=  1;  liberar(&s
>bloqueio);
 }
 O  fato  de  P  conter  s->lock  impede  que  V  tente  acordá-lo  entre  a  verificação  de  s->count  por  P  e  sua  
chamada  para  dormir.  Observe,  no  entanto,  que  precisamos  de  dormir  para  liberar  atomicamente  s->lock  e  
colocar  o  processo  consumidor  em  sono,  a  fim  de  evitar  despertares  perdidos.
 7.6  Código:  Sono  e  despertar
 Suspensão  do  Xv6  (kernel/proc.c:536)  e  despertar  (kernel/proc.c:567)  Fornece  a  interface  mostrada  no  último  
exemplo  acima,  e  sua  implementação  (além  das  regras  de  como  usá-las)  garante  que  não  haja  despertares  
perdidos.  A  ideia  básica  é  fazer  com  que  o  sleep  marque  o  processo  atual  como  SLEEPING  e,  em  seguida,  
chamar  sched  para  liberar  a  CPU;  o  wakeup  procura  um  processo  dormindo  no  canal  de  espera  fornecido  e  o  
marca  como  RUNNABLE.  Chamadores  de  sleep  e  wakeup  podem  usar  qualquer  número  mutuamente  
conveniente  como  canal.  O  Xv6  frequentemente  usa  o  endereço  de  uma  estrutura  de  dados  do  kernel  envolvida  na  espera.
 sleep  adquire  p->lock  (kernel/proc.c:547).  Agora  o  processo  que  vai  dormir  mantém  ambos  p->lock
 e  lk.  Manter  lk  era  necessário  no  chamador  (no  exemplo,  P):  isso  garantia  que  nenhum  outro  processo  (no  
exemplo,  um  V  em  execução)  pudesse  iniciar  uma  chamada  para  wakeup(chan).  Agora  que  sleep  mantém  p
>lock,  é  seguro  liberar  lk:  algum  outro  processo  pode  iniciar  uma  chamada  para  wakeup(chan),  mas  wakeup  
aguardará  para  adquirir  p->lock  e,  portanto,  aguardará  até  que  sleep  termine  de  colocar  o  processo  para  dormir,  
evitando  que  wakeup  perca  o  sleep.
 Agora  que  o  sleep  contém  p->lock  e  nenhum  outro,  ele  pode  colocar  o  processo  para  dormir  gravando  o  
canal  de  sono,  alterando  o  estado  do  processo  para  SLEEPING  e  chamando  sched  (kernel/proc.c:551-554).
 Em  breve  ficará  claro  por  que  é  essencial  que  p->lock  não  seja  liberado  (pelo  agendador)  até  que  o  processo  
seja  marcado  como  SLEEPING.
 Em  algum  momento,  um  processo  adquirirá  o  bloqueio  de  condição,  definirá  a  condição  que  o  sleeper  está  aguardando  e  
chamará  wakeup(chan).  É  importante  que  wakeup  seja  chamado  enquanto  a  condição  lock1  estiver  sendo  mantida .  O  wakeup  
percorre  a  tabela  de  processos  (kernel/proc.c:567).  Ele  adquire  o  p->lock  de  cada  processo  que  inspeciona,  tanto  porque  pode  
manipular  o  estado  desse  processo  quanto  porque  o  p->lock  garante  que  o  modo  de  suspensão  e  o  modo  de  ativação  não  se  percam.  
Quando  o  modo  de  ativação  encontra  um  processo  no  estado  SLEEPING  com  um  canal  correspondente,  ele  altera  o  estado  desse  
processo  para  RUNNABLE.  Na  próxima  vez  que  o  agendador  for  executado,  ele  verá  que  o  processo  está  pronto  para  ser  executado.
 1A  rigor,  é  suficiente  que  o  despertar  siga  apenas  a  aquisição  (ou  seja,  pode-se  chamar  o  despertar  após  a  liberação).
 78
Machine Translated by Google
 Por  que  as  regras  de  bloqueio  para  sleep  e  wakeup  garantem  que  um  processo  em  modo  sleep  não  perca  
um  wakeup?  O  processo  em  modo  sleep  mantém  o  bloqueio  de  condição  ou  seu  próprio  p->lock ,  ou  ambos,  de  
um  ponto  antes  de  verificar  a  condição  até  um  ponto  depois  de  ser  marcado  como  SLEEPING.  O  processo  que  
chama  o  wakeup  mantém  ambos  os  bloqueios  no  loop  do  wakeup.  Assim,  o  processo  que  o  ativa  torna  a  
condição  verdadeira  antes  que  a  thread  consumidora  a  verifique;  ou  o  processo  que  o  ativa  examina  a  thread  
em  modo  sleep  estritamente  após  ela  ter  sido  marcada  como  SLEEPING.  Então,  o  processo  em  modo  wakeup  
verá  o  processo  em  modo  sleep  e  o  ativará  (a  menos  que  algo  o  ative  primeiro).
 Às  vezes,  vários  processos  estão  dormindo  no  mesmo  canal;  por  exemplo,  mais  de  um  processo  lendo  de  
um  pipe.  Uma  única  chamada  para  wakeup  despertará  todos  eles.  Um  deles  será  executado  primeiro  e  obterá  o  
bloqueio  com  o  qual  sleep  foi  chamado  e  (no  caso  de  pipes)  lerá  quaisquer  dados  que  estejam  aguardando  no  
pipe.  Os  outros  processos  descobrirão  que,  apesar  de  terem  sido  despertados,  não  há  dados  para  serem  lidos.  
Do  ponto  de  vista  deles,  o  wakeup  foi  "espúrio"  e  eles  devem  dormir  novamente.  Por  esse  motivo,  sleep  é  sempre  
chamado  dentro  de  um  loop  que  verifica  a  condição.
 Não  há  problema  algum  se  dois  usuários  de  sleep/wakeup  escolherem  acidentalmente  o  mesmo  canal:  eles  
verão  despertares  espúrios,  mas  o  looping  descrito  acima  tolerará  esse  problema.  Grande  parte  do  charme  de  
sleep/wakeup  reside  em  sua  leveza  (não  há  necessidade  de  criar  estruturas  de  dados  especiais  para  atuar  como  
canais  de  sleep)  e  em  fornecer  uma  camada  de  indireção  (os  chamadores  não  precisam  saber  com  qual  processo  
específico  estão  interagindo).
 7.7  Código:  Pipes
 Um  exemplo  mais  complexo  que  usa  sleep  e  wakeup  para  sincronizar  produtores  e  consumidores  é  a  
implementação  de  pipes  do  xv6.  Vimos  a  interface  para  pipes  no  Capítulo  1:  bytes  gravados  em  uma  extremidade  
de  um  pipe  são  copiados  para  um  buffer  no  kernel  e,  em  seguida,  podem  ser  lidos  da  outra  extremidade  do  pipe.
 Os  próximos  capítulos  examinarão  o  suporte  ao  descritor  de  arquivo  em  torno  dos  pipes,  mas  vamos  agora  dar  
uma  olhada  nas  implementações  de  pipewrite  e  piperead.
 Cada  pipe  é  representado  por  um  pipe  struct,  que  contém  um  bloqueio  e  um  buffer  de  dados .
 Os  campos  nread  e  nwrite  contam  o  número  total  de  bytes  lidos  e  gravados  no  buffer.  O  buffer  é  encapsulado:  o  
próximo  byte  escrito  após  buf[PIPESIZE-1]  é  buf[0].  As  contagens  não  são  encapsuladas.  Essa  convenção  
permite  que  a  implementação  distinga  um  buffer  cheio  (nwrite  ==  nread+PIPESIZE)  de  um  buffer  vazio  (nwrite  
==  nread),  mas  significa  que  a  indexação  no  buffer  deve  usar  buf[nread  %  PIPESIZE]  em  vez  de  apenas  
buf[nread]  (e  similarmente  para  nwrite).
 Vamos  supor  que  as  chamadas  para  piperead  e  pipewrite  aconteçam  simultaneamente  em  duas  CPUs  
diferentes.  pipewrite  (kernel/pipe.c:77)  começa  adquirindo  o  bloqueio  do  pipe,  que  protege  as  contagens,  os  
dados  e  seus  invariantes  associados.  piperead  (kernel/pipe.c:106)  Então  tenta  adquirir  o  bloqueio  também,  mas  
não  consegue.  Ele  gira  em  adquirir  (kernel/spinlock.c:22)  aguardando  o  bloqueio.  Enquanto  piperead  aguarda,  
pipewrite  percorre  os  bytes  que  estão  sendo  escritos  (addr[0..n-1]),  adicionando  cada  um  ao  pipe  por  vez  (kernel/
 pipe.c:95).  Durante  esse  loop,  pode  acontecer  que  o  buffer  fique  cheio  (kernel/pipe.c:88).  Neste  caso,  o  pipewrite  
chama  wakeup  para  alertar  os  leitores  adormecidos  sobre  o  fato  de  que  há  dados  aguardando  no  buffer  e,  então,  
entra  em  modo  sleep  em  &pi->nwrite  para  esperar  que  um  leitor  retire  alguns  bytes  do  buffer.  O  sleep  libera  pi
>lock  como  parte  do  processo  de  colocar  o  processo  do  pipewrite  em  modo  sleep.
 79
Machine Translated by Google
 Agora  que  pi->lock  está  disponível,  piperead  consegue  adquiri-lo  e  entra  em  sua  seção  crítica:  ele  descobre  que  pi
>nread !=  pi->nwrite  (kernel/pipe.c:113)  (pipewrite  entrou  em  modo  de  espera  porque  pi->nwrite  ==  pi->nread+PIPESIZE  
(kernel/pipe.c:88)),  então  ele  cai  no  loop  for ,  copia  os  dados  do  pipe  (kernel/pipe.c:120),  e  incrementa  nread  pelo  número  
de  bytes  copiados.  Essa  quantidade  de  bytes  agora  está  disponível  para  escrita,  então  piperead  chama  wakeup  ( kernel/
 pipe.c:127)  para  despertar  quaisquer  escritores  adormecidos  antes  que  ele  retorne.  O  comando  wakeup  encontra  um  
processo  adormecido  em  &pi->nwrite,  o  processo  que  estava  executando  o  pipewrite,  mas  parou  quando  o  buffer  ficou  
cheio.  Ele  marca  esse  processo  como  EXECUTÁVEL.
 O  código  do  pipe  usa  canais  de  suspensão  separados  para  leitor  e  gravador  (pi->nread  e  pi->nwrite);
 Isso  pode  tornar  o  sistema  mais  eficiente  no  caso  improvável  de  haver  muitos  leitores  e  escritores  
aguardando  o  mesmo  pipe.  O  código  do  pipe  dorme  dentro  de  um  loop  que  verifica  a  condição  de  sono;  se  
houver  vários  leitores  ou  escritores,  todos  os  processos,  exceto  o  primeiro  a  acordar,  verão  que  a  condição  
ainda  é  falsa  e  dormirão  novamente.
 7.8  Código:  Espere,  saia  e  mate
 sleep  e  wakeup  podem  ser  usados  para  muitos  tipos  de  espera.  Um  exemplo  interessante,  apresentado  no  Capítulo  1,  é  
a  interação  entre  a  saída  de  um  filho  e  a  espera  de  seu  pai .  No  momento  da  morte  do  filho,  o  pai  pode  já  estar  dormindo  
em  espera  ou  fazendo  outra  coisa;  neste  último  caso,  uma  chamada  subsequente  para  wait  deve  observar  a  morte  do  
filho,  talvez  muito  tempo  depois  de  chamar  exit.  A  maneira  como  o  xv6  registra  a  morte  do  filho  até  que  wait  a  observe  é  
para  exit  colocar  o  chamador  no  estado  ZOMBIE ,  onde  permanece  até  que  wait  do  pai  perceba,  altere  o  estado  do  filho  
para  UNUSED,  copie  o  status  de  saída  do  filho  e  retorne  o  ID  do  processo  do  filho  ao  pai.  Se  o  pai  sair  antes  do  filho,  o  
pai  entrega  o  filho  ao  processo  init ,  que  chama  wait  perpetuamente;  assim,  cada  filho  tem  um  pai  para  limpar  depois  
dele.  Um  desafio  é  evitar  disputas  e  deadlocks  entre  esperas  e  saídas  simultâneas  de  pai  e  filho,  bem  como  saídas  e  
saídas  simultâneas.
 wait  começa  adquirindo  wait_lock  (kernel/proc.c:391).  O  motivo  é  que  wait_lock  atua  como  um  bloqueio  
de  condição  que  ajuda  a  garantir  que  o  pai  não  perca  um  despertar  de  um  filho  existente.
 Em  seguida,  wait  varre  a  tabela  de  processos.  Se  encontrar  um  filho  no  estado  ZOMBIE,  libera  os  recursos  desse  filho  e  
sua  estrutura  proc ,  copia  o  status  de  saída  do  filho  para  o  endereço  fornecido  para  wait  (se  não  for  0)  e  retorna  o  ID  do  
processo  do  filho.  Se  wait  encontrar  filhos,  mas  nenhum  tiver  saído,  chama  sleep  para  aguardar  que  algum  deles  saia  
(kernel/proc.c:433).  então  verifica  novamente.  wait  geralmente  contém  dois  bloqueios,  wait_lock  e  pp->lock  de  algum  
processo ;  a  ordem  para  evitar  deadlock  é  primeiro  wait_lock  e  depois  pp->lock.  exit  (kernel/proc.c:347)  Registra  o  status  
de  saída,  libera  
alguns  recursos,  chama  reparent  para  entregar  seus  filhos  ao  processo  init ,  desperta  o  pai  caso  esteja  em  espera,  
marca  o  chamador  como  zumbi  e  libera  a  CPU  permanentemente.  exit  mantém  wait_lock  e  p->lock  durante  esta  
sequência.  Ele  mantém  wait_lock  porque  é  o  bloqueio  condicional  para  o  wakeup(p->parent),  evitando  que  um  pai  em  
espera  perca  o  wakeup.  exit  também  deve  manter  p->lock  para  esta  sequência,  para  evitar  que  um  pai  em  espera  veja  
que  o  filho  está  no  estado  ZOMBIE  antes  que  o  filho  finalmente  chame  swtch.  exit  adquire  esses  bloqueios  na  mesma  
ordem  que  wait  para  evitar  deadlock.
 80
Machine Translated by Google
 Pode  parecer  incorreto  que  exit  acorde  o  pai  antes  de  definir  seu  estado  como  ZOMBIE,  mas  isso  é  seguro:  embora  
wakeup  possa  fazer  com  que  o  pai  seja  executado,  o  loop  em  wait  não  pode  examinar  o  filho  até  que  o  p->lock  do  filho  seja  
liberado  pelo  planejador,  então  wait  não  pode  examinar  o  processo  de  saída  até  bem  depois  que  exit  tenha  definido  seu  
estado  como  ZOMBIE  (kernel/proc.c:379).
 Enquanto  exit  permite  que  um  processo  se  encerre,  kill  (kernel/proc.c:586)  permite  que  um  processo  solicite  o  término  
de  outro.  Seria  muito  complexo  para  kill  destruir  diretamente  o  processo  da  vítima,  já  que  a  vítima  pode  estar  executando  em  
outra  CPU,  talvez  no  meio  de  uma  sequência  sensível  de  atualizações  nas  estruturas  de  dados  do  kernel.  Portanto,  kill  faz  
muito  pouco:  apenas  define  o  p->killed  da  vítima  e,  se  estiver  em  modo  de  espera,  a  desperta.  Eventualmente,  a  vítima  
entrará  ou  sairá  do  kernel,  momento  em  que  o  código  em  usertrap  chamará  exit  se  p->killed  estiver  definido  (ele  verifica  
chamando  killed  (kernel/proc.c:615)).  Se  a  vítima  estiver  sendo  executada  no  espaço  do  usuário,  ela  logo  entrará  no  kernel  
fazendo  uma  chamada  de  sistema  ou  porque  o  timer  (ou  algum  outro  dispositivo)  interrompe.
 Se  o  processo  da  vítima  estiver  em  modo  de  espera,  a  chamada  de  ativação  do  comando  kill  fará  com  que  a  vítima  
retorne  do  modo  de  espera.  Isso  é  potencialmente  perigoso,  pois  a  condição  que  está  sendo  aguardada  pode  não  ser  verdadeira.
 No  entanto,  chamadas  xv6  para  sleep  são  sempre  encapsuladas  em  um  loop  while  que  testa  novamente  a  condição  após  o  
retorno  do  sleep .  Algumas  chamadas  para  sleep  também  testam  p->killed  no  loop  e  abandonam  a  atividade  atual  se  ela  
estiver  definida.  Isso  só  é  feito  quando  tal  abandono  for  correto.  Por  exemplo,  o  código  de  leitura  e  gravação  do  pipe  retorna  
se  o  sinalizador  killed  estiver  definido;  eventualmente,  o  código  retornará  para  trap,  que  verificará  novamente  p->killed  e  sairá.
 Alguns  loops  de  suspensão  xv6  não  verificam  p->killed  porque  o  código  está  no  meio  de  uma  chamada  de  
sistema  multietapas  que  deveria  ser  atômica.  O  driver  virtio  (kernel/virtio_disk.c:285)  Um  exemplo:  ele  não  verifica  
p->killed  porque  uma  operação  de  disco  pode  ser  uma  de  um  conjunto  de  gravações  necessárias  para  que  o  
sistema  de  arquivos  permaneça  em  um  estado  correto.  Um  processo  que  é  encerrado  enquanto  aguarda  E/S  de  
disco  não  será  encerrado  até  concluir  a  chamada  de  sistema  atual  e  o  usertrap  visualizar  o  sinalizador  de  encerramento.
 7.9  Bloqueio  de  Processo
 O  bloqueio  associado  a  cada  processo  (p->lock)  é  o  bloqueio  mais  complexo  no  xv6.  Uma  maneira  simples  de  
pensar  em  p->lock  é  que  ele  deve  ser  mantido  durante  a  leitura  ou  gravação  de  qualquer  um  dos  seguintes  
campos  struct  proc :  p->state,  p->chan,  p->killed,  p->xstate  e  p->pid.  Esses  campos  podem  ser  usados  por  outros  
processos  ou  por  threads  do  escalonador  em  outros  núcleos,  portanto,  é  natural  que  sejam  protegidos  por  um  
bloqueio.
 Entretanto,  a  maioria  dos  usos  de  p->lock  protegem  aspectos  de  nível  superior  da  estrutura  de  dados  do  processo  xv6.
 estruturas  e  algoritmos.  Aqui  está  o  conjunto  completo  de  coisas  que  p->lock  faz:
 •  Junto  com  p->state,  ele  evita  corridas  na  alocação  de  slots  proc[]  para  novos  processos.
 •  Ele  oculta  um  processo  da  vista  enquanto  ele  está  sendo  criado  ou  destruído.
 •  Impede  que  a  espera  de  um  pai  colete  um  processo  que  definiu  seu  estado  como  ZOMBIE,  mas  tem
 ainda  não  rendeu  a  CPU.
 •  Impede  que  o  planejador  de  outro  núcleo  decida  executar  um  processo  de  produção  após  definir  seu
 estado  para  RUNNABLE ,  mas  antes  de  terminar,  alterne.
 81
Machine Translated by Google
 •  Garante  que  apenas  o  agendador  de  um  núcleo  decida  executar  um  processo  RUNNABLE .
 •  Evita  que  uma  interrupção  do  temporizador  faça  com  que  um  processo  ceda  enquanto  estiver  em  comutação.
 •  Junto  com  o  bloqueio  de  condição,  ajuda  a  evitar  que  o  despertar  ignore  um  processo  que  está
 chamando  sleep  mas  não  terminou  de  renderizar  a  CPU.
 •  Impede  que  o  processo  de  morte  da  vítima  saia  e  talvez  seja  realocado  entre
 verificação  de  p->pid  e  configuração  de  p->killed.
 •  Torna  a  verificação  e  gravação  de  kill  de  p->state  atômicas.
 O  campo  p->parent  é  protegido  pelo  bloqueio  global  wait_lock  em  vez  de  p->lock.
 Somente  o  processo  pai  modifica  p->parent,  embora  o  campo  seja  lido  tanto  pelo  próprio  processo  quanto  por  
outros  processos  que  buscam  seus  filhos.  O  propósito  de  wait_lock  é  atuar  como  o  bloqueio  condicional  quando  
wait  está  em  modo  sleep,  aguardando  a  saída  de  qualquer  filho.  Um  filho  que  está  saindo  mantém  wait_lock  ou  p
>lock  até  que  seu  estado  seja  definido  como  ZOMBIE,  seu  pai  tenha  despertado  e  a  CPU  tenha  sido  liberada.  
wait_lock  também  serializa  saídas  simultâneas  de  um  pai  e  um  filho,  de  modo  que  o  processo  init  (que  herda  o  
filho)  tenha  a  garantia  de  ser  despertado  de  sua  espera.  wait_lock  é  um  bloqueio  global,  e  não  um  bloqueio  por  
processo  em  cada  pai,  porque,  até  que  um  processo  o  adquira,  ele  não  pode  saber  quem  é  seu  pai.
 7.10  Mundo  real
 O  escalonador  xv6  implementa  uma  política  de  escalonamento  simples,  que  executa  cada  processo  por  vez.  Essa  
política  é  chamada  de  round  robin.  Sistemas  operacionais  reais  implementam  políticas  mais  sofisticadas  que,  por  
exemplo,  permitem  que  os  processos  tenham  prioridades.  A  ideia  é  que  um  processo  executável  de  alta  prioridade  
seja  preferido  pelo  escalonador  a  um  processo  executável  de  baixa  prioridade.  Essas  políticas  podem  se  tornar  
complexas  rapidamente  porque  frequentemente  há  objetivos  conflitantes:  por  exemplo,  o  sistema  operacional  
também  pode  querer  garantir  justiça  e  alto  rendimento.  Além  disso,  políticas  complexas  podem  levar  a  interações  
não  intencionais,  como  inversão  de  prioridade  e  comboios.  A  inversão  de  prioridade  pode  ocorrer  quando  um  
processo  de  baixa  prioridade  e  um  de  alta  prioridade  usam  um  bloqueio  específico,  que,  quando  adquirido  pelo  
processo  de  baixa  prioridade,  pode  impedir  o  processo  de  alta  prioridade  de  progredir.  Um  longo  comboio  de  
processos  em  espera  pode  se  formar  quando  muitos  processos  de  alta  prioridade  aguardam  por  um  processo  de  
baixa  prioridade  que  adquire  um  bloqueio  compartilhado;  uma  vez  formado,  o  comboio  pode  persistir  por  um  longo  
tempo.  Para  evitar  esses  tipos  de  problemas,  mecanismos  adicionais  são  necessários  em  planejadores  sofisticados.
 Sleep  e  Wakeup  são  métodos  de  sincronização  simples  e  eficazes,  mas  existem  muitos  outros.  O  primeiro  
desafio  em  todos  eles  é  evitar  o  problema  de  "perdas  de  ativação"  que  vimos  no  início  do  capítulo.  A  função  sleep  
do  kernel  Unix  original  simplesmente  desabilitava  interrupções,  o  que  era  suficiente  porque  o  Unix  rodava  em  um  
sistema  com  uma  única  CPU.  Como  o  xv6  roda  em  multiprocessadores,  ele  adiciona  um  bloqueio  explícito  à  
função  sleep.  O  msleep  do  FreeBSD  adota  a  mesma  abordagem.  A  função  sleep  do  Plan  9  usa  uma  função  de  
retorno  de  chamada  que  é  executada  com  o  bloqueio  de  agendamento  mantido  pouco  antes  de  entrar  em  modo  
sleep;  a  função  serve  como  uma  verificação  de  última  hora  da  condição  de  sleep,  para  evitar  perdas  de  ativação.  O  kernel  Linux
 82
Machine Translated by Google
 sleep  usa  uma  fila  de  processos  explícita,  chamada  fila  de  espera,  em  vez  de  um  canal  de  espera;  a  fila  tem  seu  
próprio  bloqueio  interno.
 Varrer  todo  o  conjunto  de  processos  durante  o  processo  de  despertar  é  ineficiente.  Uma  solução  melhor  é  substituir  o  canal  
tanto  no  modo  de  suspensão  quanto  no  de  despertar  por  uma  estrutura  de  dados  que  contenha  uma  lista  de  processos  em  
suspensão  nessa  estrutura,  como  a  fila  de  espera  do  Linux.  O  modo  de  suspensão  e  o  modo  de  despertar  do  Plan  9  chamam  essa  
estrutura  de  ponto  de  encontro.  Muitas  bibliotecas  de  threads  se  referem  à  mesma  estrutura  como  uma  variável  de  condição;  
nesse  contexto,  as  operações  de  suspensão  e  despertar  são  chamadas  de  espera  e  sinalização.  Todos  esses  mecanismos  
compartilham  o  mesmo  princípio:  a  condição  de  suspensão  é  protegida  por  algum  tipo  de  bloqueio  que  é  descartado  atomicamente  
durante  o  processo  de  suspensão.
 A  implementação  do  wakeup  ativa  todos  os  processos  que  aguardam  em  um  canal  específico,  e  pode  ser  que  muitos  
processos  estejam  aguardando  esse  canal  específico.  O  sistema  operacional  agendará  todos  esses  processos  e  eles  verificarão  
rapidamente  a  condição  de  suspensão.
 Processos  que  se  comportam  dessa  maneira  são  às  vezes  chamados  de  "manada  estrondosa"  e  é  melhor  evitá-los.
 A  maioria  das  variáveis  de  condição  tem  duas  primitivas  para  ativação:  sinal,  que  ativa  um  processo,  e  transmissão,  
que  ativa  todos  os  processos  em  espera.
 Semáforos  são  frequentemente  usados  para  sincronização.  A  contagem  normalmente  corresponde  a  algo  como  o  número  de  
bytes  disponíveis  em  um  buffer  de  pipe  ou  o  número  de  filhos  zumbis  que  um  processo  possui.  Usar  uma  contagem  explícita  como  
parte  da  abstração  evita  o  problema  do  "despertar  perdido":  há  uma  contagem  explícita  do  número  de  despertares  ocorridos.  A  
contagem  também  evita  os  problemas  de  despertares  espúrios  e  de  manada  estrondosa.
 Terminar  processos  e  limpá-los  introduz  muita  complexidade  no  xv6.  Na  maioria  dos  sistemas  operacionais,  é  ainda  
mais  complexo,  porque,  por  exemplo,  o  processo  vítima  pode  estar  dormindo  profundamente  dentro  do  kernel,  e  desenrolar  
sua  pilha  requer  cuidado,  já  que  cada  função  na  pilha  de  chamadas  pode  precisar  fazer  alguma  limpeza.  Algumas  
linguagens  ajudam  fornecendo  um  mecanismo  de  exceção,  mas  não  C.  Além  disso,  existem  outros  eventos  que  podem  
fazer  com  que  um  processo  em  espera  seja  despertado,  mesmo  que  o  evento  que  ele  está  aguardando  ainda  não  tenha  
ocorrido.  Por  exemplo,  quando  um  processo  Unix  está  em  espera,  outro  processo  pode  enviar  um  sinal  para  ele.  Nesse  
caso,  o  processo  retornará  da  chamada  de  sistema  interrompida  com  o  valor  -1  e  com  o  código  de  erro  definido  como  
EINTR.  O  aplicativo  pode  verificar  esses  valores  e  decidir  o  que  fazer.  O  Xv6  não  suporta  sinais  e  essa  complexidade  não  
surge.
 O  suporte  do  Xv6  para  kill  não  é  totalmente  satisfatório:  existem  loops  de  sleep  que  provavelmente  deveriam  verificar  se  há  p
>killed.  Um  problema  relacionado  é  que,  mesmo  para  loops  de  sleep  que  verificam  p->killed,  há  uma  corrida  entre  sleep  e  kill;  este  
último  pode  ativar  p->killed  e  tentar  acordar  a  vítima  logo  após  o  loop  da  vítima  verificar  p->killed,  mas  antes  de  chamar  sleep.  Se  
esse  problema  ocorrer,  a  vítima  não  notará  o  p->killed  até  que  a  condição  que  está  aguardando  ocorra.  Isso  pode  ocorrer  um  
pouco  mais  tarde  ou  até  mesmo  nunca  (por  exemplo,  se  a  vítima  estiver  aguardando  uma  entrada  do  console,  mas  o  usuário  não  
digitar  nenhuma  entrada).
 Um  sistema  operacional  real  encontraria  estruturas  proc  livres  com  uma  lista  livre  explícita  em  constante
 tempo  em  vez  da  pesquisa  de  tempo  linear  em  allocproc;  xv6  usa  a  varredura  linear  para  simplificar.
